\section{Introducción}

El objetivo del proyecto es programar en assembly un juego de reflejos que corra sobre
un procesador RISC-V real, instanciado en una FPGA, y traducir ese programa a código de
máquina con un assembler desarrollado por el grupo. El punto de partida es el Pochoco
SoC \cite{pochoco}: un sistema en chip con el procesador Espino (RV32E), una memoria RAM
unificada para instrucciones y datos, y periféricos mapeados en memoria para los LEDs,
los botones y los displays de la Go Board \cite{goboard}.

El trabajo tiene tres partes, que corresponden a las secciones de este informe:

\begin{enumerate}
  \item \textbf{Hardware} (sección~\ref{sec:arquitectura}). El SoC original no ofrece
        ninguna forma de medir tiempo. Se agregó un contador de ciclos de reloj que el
        programa puede leer, y un módulo superior \code{game\_top} que instancia el SoC
        con la imagen de memoria del juego.
  \item \textbf{Software} (secciones~\ref{sec:temporizacion} y~\ref{sec:programa}). El
        juego en \code{sw/game.s}: temporización de tres segundos, medición del tiempo
        de respuesta, conversión de ciclos a décimas de segundo, promedio de las diez
        rondas y elección pseudoaleatoria del LED objetivo.
  \item \textbf{Assembler} (sección~\ref{sec:assembler}). El programa
        \code{asm/pochoasm.py}, que traduce \code{game.s} a \code{game.hex}: una palabra
        de 32 bits en hexadecimal por línea, que la RAM del SoC carga al sintetizar el
        diseño.
\end{enumerate}

La sección~\ref{sec:verificacion} describe cómo se verificó cada parte, y la
sección~\ref{sec:operacion} cómo compilar, programar la placa y cambiar una constante
del juego, que es lo que se pedirá durante la demostración.

\subsection{Reglas del juego}
\label{sec:reglas}

La secuencia implementada sigue el enunciado:

\begin{enumerate}
  \item Al encender, los LEDs parpadean alternando 0101 y 1010 y los displays muestran
        \code{00}. El juego espera que se apriete cualquier botón para empezar.
  \item Cada ronda comienza con los cuatro LEDs encendidos durante al menos tres
        segundos.
  \item Luego se enciende solo el LED objetivo, elegido de forma pseudoaleatoria. En
        ese mismo instante comienza la medición de tiempo.
  \item El jugador aprieta un botón. Si es el del LED objetivo, los displays muestran
        el tiempo de respuesta en décimas de segundo y la ronda cuenta como correcta.
  \item Si es otro botón, los displays muestran \code{EE} y la ronda se repite desde
        el paso~2. Las rondas correctas acumuladas se conservan.
  \item Tras diez respuestas correctas, el último tiempo queda visible tres segundos
        y después se muestra el promedio, con los LEDs parpadeando. Un botón inicia
        una partida nueva.
\end{enumerate}

Sobre el paso~3, el enunciado pide iniciar la medición ``en el mismo instante lógico en
que comienza la ronda''. Se interpretó que el tiempo de respuesta se cuenta desde que
aparece el estímulo, es decir, desde que se enciende el LED objetivo: si se contara
desde que se encienden los cuatro LEDs, todos los tiempos incluirían los tres segundos
fijos de la primera fase. La lectura del contador que inicia la medición es la
instrucción inmediatamente siguiente a la que enciende el LED objetivo.
